% Multimodal MIL — Lung Transplant Outcome Prediction with Interpretability
% Compile with: pdflatex model_description.tex
\documentclass[11pt,a4paper]{article}

\usepackage{amsmath,amssymb,amsfonts}
\usepackage{booktabs}
\usepackage{graphicx}
\usepackage{xcolor}
\usepackage{hyperref}
\usepackage{geometry}
\usepackage{microtype}
\usepackage[numbers,sort&compress]{natbib}
\geometry{margin=2.5cm}

\title{\textbf{Multimodal Multiple Instance Learning\\
for Lung Transplant Outcome Prediction\\
with Interpretability}}
\author{Dinesh Haridoss}
\date{\today}

\begin{document}
\maketitle
\tableofcontents
\newpage

% =============================================================================
\section{Overview}
% =============================================================================

We present a two-phase multimodal Multiple Instance Learning (MIL) framework for
classifying Acute Cellular Rejection (ACR) in lung transplant patients.
The framework ingests up to four modalities: H\&E histology (HE), BAL cytology
(BAL), computed tomography (CT), and clinical tabular data (Clinical).
Each modality is represented as a \emph{bag of instances} without requiring
spatial registration across modalities.

H\&E patches carry spatial tile coordinates (\texttt{tile\_left},
\texttt{tile\_top}) that record each patch's pixel position on the biopsy slide.
This spatial information enables position-aware visualisation of attention maps
on the biopsy tissue.
CT instances carry 3-D voxel coordinates (\texttt{d\_voxel}, \texttt{h\_voxel},
\texttt{w\_voxel}) locating each patch within the scan volume.

\paragraph{Phase 1.}
Each modality is trained independently using a Gated Attention MIL
backbone~\cite{ilse2018abmil}.
Optional auxiliary objectives encourage student modalities (CT, BAL) to mimic
the representations of teacher modalities (HE, Clinical).
These objectives include cross-modal attention transfer, contrastive
representation alignment, and cosine representation distillation.

\paragraph{Phase 2.}
Pre-trained modality encoders are frozen (projection heads) or fine-tuned
(backbones), and a fusion module combines all available modality representations.
We evaluate five fusion strategies: early, late, middle, cross-attention, and
iterative cross-modal fusion.

\bigskip\noindent
\textbf{Task.}
Binary classification (rejection vs.\ non-rejection) under extreme class
imbalance.
The model is evaluated by balanced accuracy (bAcc) and AUROC.

% =============================================================================
\section{Dataset Preprocessing}
% =============================================================================

\subsection{Raw Data Sources}

The dataset is assembled from five raw sources covering four imaging and one
tabular modality:

\begin{table}[h]
\centering
\caption{Raw data sources and their identifiers.}
\label{tab:datasources}
\begin{tabular}{llll}
\toprule
Modality & File type & Patient ID col & Time col \\
\midrule
BAL (scRNA-seq) & \texttt{.h5ad} & \texttt{record\_id} & \texttt{date\_from\_id} \\
H\&E histology  & \texttt{.h5ad} & \texttt{record\_id} & \texttt{biopsy\_date} \\
CT scan         & \texttt{.h5ad} & \texttt{Patient}    & \texttt{date\_of\_CT} \\
Radiomics       & \texttt{.csv}  & \texttt{Patient}    & \texttt{Exam\_Date} \\
Clinical        & \texttt{.csv}  & \texttt{record\_id} & \texttt{spiro\_date} \\
\bottomrule
\end{tabular}
\end{table}

\noindent
Instance-level embeddings for the three imaging modalities are stored in the
\texttt{.obsm} key \texttt{X\_scVI} (BAL) or \texttt{X} (H\&E, CT).
Per-instance cluster and cell-type labels are stored in the \texttt{.obs}
columns \texttt{resolution\_v2} (BAL), \texttt{subcluster\_renamed} (H\&E),
and \texttt{leiden} (CT).
For H\&E, spatial coordinates \texttt{tile\_left} and \texttt{tile\_top}
record the pixel position of each patch on the biopsy slide.
For CT, coordinates \texttt{d\_voxel}, \texttt{h\_voxel}, and
\texttt{w\_voxel} record the 3-D voxel location of each patch in the scan
volume.

\subsection{Dataset Statistics and Modality Coverage}

The cohort consists of longitudinal biopsy-anchored timepoints from lung
transplant patients.
Each timepoint is associated with a rejection label derived from the ACR
biopsy grade.
Figure~\ref{fig:modality_overlap} shows the patient-level modality availability
across the cohort.
The largest single-modality group consists of Clinical-only patients (257
patients), reflecting that clinical measurements are the most frequently
collected.
A total of 60 patients have all four modalities available simultaneously.

Figure~\ref{fig:acr_distribution} shows the distribution of ACR-positive and
ACR-negative timepoints for each modality combination.
The dataset is highly imbalanced: rejection events are a minority across all
modality groups.
Clinical-only timepoints are the most frequent (1643 ACR$-$, 425 ACR$+$),
and the HE+CT+Clinical group contributes 169 ACR$-$ and 64 ACR$+$ timepoints.

\begin{figure}[h]
\centering
\includegraphics[width=0.85\textwidth]{../chicago/plots/Patient-Level_modality_overlap.png}
\caption{Patient-level modality overlap (UpSet plot). Each bar represents a
  modality combination. The largest single-modality group is Clinical-only
  (257 patients). The all-four-modality group contains 60 patients.}
\label{fig:modality_overlap}
\end{figure}

\begin{figure}[h]
\centering
\includegraphics[width=0.85\textwidth]{../chicago/plots/Timepoint-Level_acr_distribution_by_modality.png}
\caption{Timepoint-level ACR label distribution by modality combination.
  Each bar shows the count of ACR-positive and ACR-negative timepoints for a
  given modality subset. Clinical-only timepoints are the most frequent.}
\label{fig:acr_distribution}
\end{figure}

\subsection{Temporal Alignment}
\label{sec:temporal}

All modalities are aligned to a common anchor event using a time-windowed
nearest-neighbour strategy.
The H\&E biopsy is the anchor modality because it provides the ACR rejection
label.
For each biopsy event, the dataset searches for the nearest available
measurement from every other modality within a window of $\pm$45 days
($45 \times 86\,400$ seconds).
If multiple measurements from the same modality fall within the window, the one
closest in time to the biopsy date is selected.
A timepoint is included as long as at least one modality is present
(\texttt{min\_modalities}~$= 1$).

The alignment procedure for each biopsy anchor proceeds as follows:
\begin{enumerate}
    \item Fix the anchor timestamp $t_0$ from the biopsy date.
    \item For each other modality $m \in \{\text{BAL, CT, Clinical}\}$, collect
          all measurements from the same patient.
    \item Retain measurements whose timestamp $t_m$ satisfies
          $|t_m - t_0| \leq 45 \times 86\,400$.
    \item Among retained measurements, select the one minimising $|t_m - t_0|$
          as the aligned record for modality $m$.
    \item If no measurement falls within the window, modality $m$ is absent for
          this timepoint.
\end{enumerate}

This procedure produces a set of aligned multimodal samples.
Each sample has one rejection label from the biopsy and zero or more
co-occurring modality observations within the 45-day window.
Patients without a biopsy on a given date still contribute if their most
recent biopsy label falls within the window.

\subsection{Label Encoding}

The binary rejection label is derived from the ACR biopsy grade:
\begin{equation}
    y =
    \begin{cases}
        1 & \text{if ACR grade contains A1 or A2 (acute rejection),} \\
        0 & \text{otherwise.}
    \end{cases}
\end{equation}
Outcome-predictive clinical columns (FVC/FEV1 slopes, biopsy grade itself) are
excluded from the clinical feature set to prevent label leakage.

\subsection{Clinical Feature Tokenisation}

Raw clinical tabular features ($F = 102$ features after label-leakage
exclusion) are discretised by a \texttt{ClinicalFeatureTokenizer}:

\paragraph{Continuous features.}
Each continuous feature $f$ is discretised into $n_q = 4$ quantile bins fitted
on the population.
Bin boundaries $b_0^{(f)} < b_1^{(f)} < \cdots < b_{n_q}^{(f)}$ are the
$0, 1/n_q, \ldots, 1$ quantiles of the observed values.
A dedicated NaN token is reserved per feature.

\paragraph{Categorical/binary features.}
Features with $\leq 10$ unique values are treated as categorical: each unique
category receives a distinct token, plus one NaN token.

\paragraph{Block-diagonal one-hot representation.}
At precompute time, each sample's tokenised clinical vector $\mathbf{t} \in
\mathbb{Z}^F$ (one token ID per feature) is converted into the 2-D bag matrix
$\mathbf{X}_\text{Clin} \in \{0,1\}^{F \times V_\text{Clin}}$.
The global vocabulary has $V_\text{Clin}$ tokens, partitioned into $F$
non-overlapping blocks of size $n_q + 1$ (continuous) or
$|\text{categories}| + 1$ (categorical).
Row $f$ of $\mathbf{X}_\text{Clin}$ is a one-hot vector with a single $1$ at
the position of token $t_f$.
Because each feature occupies a disjoint block of columns, the matrix is
block-diagonal and each row is a valid token embedding.
This treats clinical features as $F$ \emph{instance tokens} on equal footing
with imaging patch tokens.

For the special case of $n_q = 4$ bins with all $F$ features continuous, the
representation simplifies to $\mathbf{X}_\text{Clin} \in \{0,1\}^{F \times 4F}$
(feature dimension $= 4F = 408$), matching the \texttt{MODALITY\_REGISTRY}
entry \texttt{("clinical\_onehot", 408)}.

\subsection{Cluster Count Tokenisation}

For each imaging modality $m \in \{\text{HE, BAL, CT}\}$ with $K_m$
cell/patch clusters, a per-bag cluster abundance vector
$\mathbf{c}_m \in \mathbb{R}^{K_m}$ (CLR-normalised counts) is precomputed
from the per-instance cluster labels.
A \texttt{ClusterCountTokenizer} discretises each cluster's abundance into
$n_q = 4$ quantile bins.
Each cluster receives its own separate set of quantile boundaries fitted on the
population distribution for that cluster.
The boundaries for cluster $k$ are computed from the CLR-normalised count of
cluster $k$ across all patients.
This per-cluster fitting ensures that the bin thresholds are meaningful for
each cell type independently, regardless of how common or rare that cell type
is.

The token ID vector is $\mathbf{t}_m \in \mathbb{Z}^{K_m}$, where entry $k$
encodes the quantile bin of cluster $k$'s abundance relative to the
population.
The analogous block-diagonal one-hot matrix is
$\mathbf{X}_m^\text{count} \in \{0,1\}^{K_m \times (n_q K_m)}$: row $k$ has a
$1$ at column $k n_q + b_k$, where $b_k$ is the quantile bin of cluster $k$'s
count.
Pooling across rows gives the bag-level count representation:
\begin{equation}
    \mathbf{c}_m^\text{onehot} = (\mathbf{X}_m^\text{count})^\top\mathbf{1}_{K_m}
    \in \{0,1\}^{n_q K_m},
\end{equation}
a sparse vector with exactly $K_m$ non-zero entries (one per cluster).

\subsection{Precomputed \texttt{.pt} File Format}

Each aligned timepoint is saved as a PyTorch \texttt{.pt} dictionary with the
following top-level keys after key remapping (\texttt{KEY\_REMAP}):

\begin{itemize}
    \item \texttt{inputs}: dict of modality tensors.
          \texttt{HE\_cells}~$(N_\text{HE}, 1024)$,
          \texttt{BAL\_cells}~$(N_\text{BAL}, 10)$,
          \texttt{CT\_cells}~$(N_\text{CT}, 1024)$,
          \texttt{Clinical}~$(N_\text{Clin}, d_\text{Clin})$.
    \item \texttt{cluster\_labels}: per-instance cluster assignment strings.
    \item \texttt{coords}: per-instance spatial coordinates.
    \item \texttt{bag\_centroids}: dict of per-modality centroid
          matrices~$(K_m, d_m)$.
    \item \texttt{clinical\_onehot}: $\mathbf{X}_\text{Clin}
          \in \{0,1\}^{F \times 4F}$.
    \item \texttt{cluster\_count\_onehot}: dict of $\mathbf{X}_m^\text{count}$.
    \item \texttt{cluster\_names}: dict of cluster name lists per modality.
    \item \texttt{clinical\_feature\_names}: list of $F$ feature name strings.
    \item \texttt{clinical\_token\_ids}: per-patient bin assignment indices.
    \item \texttt{clinical\_vocab}: list of bin-specific token label dicts
          (e.g.\ \texttt{\{"id": 3, "label": "FEV1\_q3"\}}).
    \item \texttt{label}, \texttt{identifier}, \texttt{anchor\_time},
          \texttt{metadata}, \texttt{survival}.
\end{itemize}

% =============================================================================
\section{Data Representation}
% =============================================================================

\subsection{Modality Feature Spaces}

\begin{table}[h]
\centering
\caption{Modality registry. Each bag is a set of instance-level feature vectors.}
\label{tab:modalities}
\begin{tabular}{llrl}
\toprule
Modality & Key & Feature dim $d$ & Role \\
\midrule
HE       & \texttt{HE\_cells}        & 1024 & Teacher \\
Clinical & \texttt{clinical\_onehot} & 408  & Teacher \\
CT       & \texttt{CT\_cells}        & 1024 & Student \\
BAL      & \texttt{BAL\_cells}       & 10   & Student \\
\bottomrule
\end{tabular}
\end{table}

\paragraph{Clinical representation.}
Clinical features are encoded as a 2-D bag
$\mathbf{X}_\text{Clin}\in\mathbb{R}^{F\times(F\cdot n_b)}$, where $F=102$ is
the number of clinical features and $n_b=4$ is the number of quantisation bins.
Row $f$ is a block-diagonal one-hot vector that encodes the binned value of
feature $f$.
There is exactly one non-zero entry at column $f\cdot n_b + b_f$, where $b_f$
is the bin index.
This representation treats each clinical feature as a separate
\emph{token instance}, enabling attention-based pooling over clinical features
on equal footing with histological patch tokens.

% =============================================================================
\section{Phase 1: Unimodal Pre-Training}
% =============================================================================

\subsection{Gated Attention MIL Encoder}
\label{sec:abmil}

Following~\cite{ilse2018abmil}, each modality encoder maps a bag
$\mathbf{X}\in\mathbb{R}^{N\times d}$ (with $N$ instances) to a scalar logit
via three stages.

\paragraph{Instance projection (backbone).}
\begin{equation}
    \mathbf{h}_i = \text{ReLU}\!\left(\mathbf{W}_\text{bb}\,\mathbf{x}_i\right),
    \quad \mathbf{h}_i \in \mathbb{R}^H, \quad H=256.
\end{equation}

\paragraph{Gated attention pooling.}
\begin{equation}
    a_i = \frac{\exp\!\left(\mathbf{w}^\top
          \!\left[\tanh(\mathbf{V}\mathbf{h}_i)
                 \odot\sigma(\mathbf{U}\mathbf{h}_i)\right]\right)}
         {\sum_{j=1}^{N}\exp\!\left(\mathbf{w}^\top
          \!\left[\tanh(\mathbf{V}\mathbf{h}_j)
                 \odot\sigma(\mathbf{U}\mathbf{h}_j)\right]\right)},
    \label{eq:abmil}
\end{equation}
where $\mathbf{V},\mathbf{U}\in\mathbb{R}^{H\times H}$,
$\mathbf{w}\in\mathbb{R}^H$.

\paragraph{Bag representation and classification.}
\begin{equation}
    \mathbf{r} = \sum_{i=1}^{N} a_i\,\mathbf{h}_i, \qquad
    \hat{y} = \mathbf{w}_\text{head}^\top \mathbf{r}.
    \label{eq:rep}
\end{equation}

\subsection{Phase 1 Loss Functions}

\subsubsection{Task Loss (Weighted Hinge)}

\begin{equation}
    \mathcal{L}_\text{task} = \frac{1}{B}\sum_{b=1}^{B}
        w_b\cdot\max\!\left(0,\; 1 - \tilde{y}_b\,\hat{y}_b\right),
    \label{eq:hinge}
\end{equation}
where $\tilde{y}_b = 2y_b - 1 \in\{-1,+1\}$ and $w_b$ is a class weight that
upweights the minority (rejection) class.

\subsubsection{Bidirectional Patch Cross-Attention}

For student--teacher pairs we compute bidirectional patch-level cross-attention
following the MCAT architecture~\cite{chen2021mcat}.
Let $\mathbf{H}_s\in\mathbb{R}^{N_s\times H}$ and
$\mathbf{H}_t\in\mathbb{R}^{N_t\times H}$ be the student and teacher patch
features.

\paragraph{Student attends teacher.}
\begin{equation}
    \text{CA}_{s\to t} = \text{softmax}\!\left(
        \frac{\mathbf{Q}_s\,\mathbf{K}_t^\top}{\sqrt{H/n_h}}\right)\mathbf{V}_t,
\end{equation}
producing enriched student features $\tilde{\mathbf{H}}_s$.
The multi-head variant uses $n_h = 4$ heads.
After cross-attention, layer normalisation and a feed-forward residual are
applied.

\paragraph{Knowledge Distillation signal.}
For knowledge distillation, the \emph{teacher attends to student embeddings}.
The teacher patches act as queries and the student patches act as keys and
values.
This produces cross-attention weights
$\boldsymbol{\alpha}_{t\to s}\in\mathbb{R}^{N_t\times N_s}$, where entry
$[j, i]$ measures how strongly teacher patch $j$ attends to student patch $i$.
These attention weights are computed from the pre-softmax logits
$\mathbf{Q}_t\mathbf{K}_s^\top / \sqrt{H/n_h}$.
The per-student-patch importance score is obtained by averaging over teacher
patches:
\begin{equation}
    \bar{\boldsymbol{\alpha}}_{t\to s} = \frac{1}{N_t}
        \sum_{j=1}^{N_t}\boldsymbol{\alpha}_{t\to s}[j,:]\;\in\mathbb{R}^{N_s}.
\end{equation}
This score tells us which student patches the teacher considers most
informative, and it serves as a soft target for the student's own
self-attention.

\subsubsection{Attention Transfer Loss (KD)}

Inspired by~\cite{zagoruyko2017at}, we align the student's self-attention
distribution to the teacher-guided distribution over the top-$K$ most attended
student patches.
Both $\boldsymbol{\alpha}_\text{self}$ and $\bar{\boldsymbol{\alpha}}_{t\to s}$
are pre-softmax attention logits.
Softmax is applied inside the KL divergence to convert them to distributions:
\begin{equation}
    \mathcal{L}_\text{KD} =
    \mathrm{KL}\!\left(
        \text{softmax}\!\left(\frac{\boldsymbol{\alpha}_\text{self}[\mathcal{K}]}{\tau}\right)
        \,\Big\|\,
        \text{softmax}\!\left(\frac{\bar{\boldsymbol{\alpha}}_{t\to s}[\mathcal{K}]}{\tau}\right)
    \right),
    \label{eq:kd}
\end{equation}
where $\mathcal{K}$ is the index set of the top-$K=50$ student patches by
teacher score, and $\tau=0.5$ is a temperature parameter.
Using pre-softmax logits allows the temperature $\tau$ to control the
sharpness of the target distribution independently of the attention scores.

\subsubsection{Cosine Representation Distillation (CRD)}

To align the student's cross-attended bag representation
$\mathbf{r}_s^\times$ with the cached teacher summary $\mathbf{r}_t$:
\begin{equation}
    \mathcal{L}_\text{CRD} = 1 - \frac{\mathbf{r}_s^\times}{\|\mathbf{r}_s^\times\|}
        \cdot\frac{\mathbf{r}_t}{\|\mathbf{r}_t\|},
    \label{eq:crd}
\end{equation}
adapted from CRD~\cite{tian2020crd}.

\subsubsection{Supervised Contrastive Loss (SupCon)}

All modality representations $\mathbf{z} = g(\mathbf{r})$ (projected through a
two-layer MLP projection head $g:\mathbb{R}^H\to\mathbb{R}^{128}$) contribute
to a SupCon loss~\cite{khosla2020supcon}:
\begin{equation}
    \mathcal{L}_\text{CLR} =
    -\sum_{i\in I} w_i
    \frac{1}{|P(i)|}\sum_{p\in P(i)}
    \log\frac{\exp(\mathbf{z}_i\cdot\mathbf{z}_p/\tau_c)}
             {\sum_{j\in I\setminus\{i\}}\exp(\mathbf{z}_i\cdot\mathbf{z}_j/\tau_c)},
    \label{eq:supcon}
\end{equation}
with $\tau_c = 0.07$ and positive set
$P(i) = \{j : (\text{same stem} \wedge \text{different modality}) \vee
             (\text{same label} \wedge \text{different stem})\}$.
Class weights $w_i$ upweight the minority class.
An analogous \emph{summary CLR} loss is computed on the normalised bag
representations $\ell_2(\mathbf{r})$ directly, without a projection head.

\subsubsection{Total Phase 1 Objective}

\begin{equation}
    \mathcal{L}_1 = \mathcal{L}_\text{task}
        + \lambda_\text{cx}\,\mathcal{L}_\text{cross}
        + \lambda_\text{kd}\,\mathcal{L}_\text{KD}
        + \lambda_\text{crd}\,\mathcal{L}_\text{CRD}
        + \lambda_\text{clr}\,\mathcal{L}_\text{CLR}
        + \lambda_\text{sclr}\,\mathcal{L}_\text{sCLR},
    \label{eq:p1_total}
\end{equation}
where $\mathcal{L}_\text{cross}$ is an auxiliary hinge loss computed from the
cross-attended student representation.
The default configuration uses task loss only, with all $\lambda$s set to 0
unless explicitly enabled.

\subsection{Phase 1 Hyperparameters}

\begin{table}[h]
\centering
\caption{Phase 1 hyperparameters.}
\label{tab:p1_hparams}
\begin{tabular}{llr}
\toprule
Parameter & Symbol & Value \\
\midrule
Hidden dimension & $H$ & 256 \\
Dropout & $p$ & 0.4 \\
Learning rate & $\eta_1$ & $1\times10^{-5}$ \\
Weight decay & & $1\times10^{-3}$ \\
Max epochs & & 400 \\
Eval every & & 25 epochs \\
Grad accumulation steps & & 4 \\
Cross-attn heads & $n_h$ & 4 \\
Projection dim & $d_z$ & 128 \\
KD temperature & $\tau$ & 0.5 \\
KD top-K & $K$ & 50 \\
CLR temperature & $\tau_c$ & 0.07 \\
$\lambda_\text{cx}$ & & 0.05 \\
$\lambda_\text{kd}$ & & 0.05 \\
$\lambda_\text{crd}$ & & 0.10 \\
$\lambda_\text{clr}$ & & 0.10 \\
$\lambda_\text{sclr}$ & & 0.10 \\
\bottomrule
\end{tabular}
\end{table}

% =============================================================================
\section{Phase 2: Multimodal Fusion}
% =============================================================================

\subsection{Shared Building Blocks}

\subsubsection{Iterative Slot Attention}
\label{sec:slots}

Slot Attention~\cite{locatello2020slots} compresses a variable-length patch
sequence $\mathbf{H}\in\mathbb{R}^{N\times H}$ into $K$ fixed-size slot
vectors $\mathbf{S}\in\mathbb{R}^{K\times H}$, where $K \ll N$.
Each slot learns to represent a different subset of the input patches through
a competition mechanism.

\paragraph{Set Transformer view.}
The core idea is related to the Induced Set Attention Block (ISAB) from the
Set Transformer~\cite{lee2019settransformer}.
Two attention operations define the block.
First, the $K$ slot vectors act as queries and the $N$ patch features act as
keys and values:
\begin{equation}
    \mathbf{S}' = \mathrm{MHA}(\mathbf{S},\,\mathbf{H},\,\mathbf{H}).
    \label{eq:mab}
\end{equation}
Each slot gathers information from the patches it is most similar to.
Second, the patches can be refined by attending back to the updated slots:
\begin{equation}
    \mathbf{H}' = \mathrm{MHA}(\mathbf{H},\,\mathbf{S}',\,\mathbf{S}').
    \label{eq:sab}
\end{equation}
Here $\mathbf{S}$ is a small set of $K$ vectors in $\mathbb{R}^d$, with
$K \ll N$.
This bottleneck forces the slots to form compact summaries of the full patch
sequence.

\paragraph{Slot Attention procedure.}
Our implementation follows~\cite{locatello2020slots} and uses a GRU-based
iterative update.
A slot-competition softmax is applied column-wise over patches so that each
patch is drawn toward the slot it fits best.

Given patch features $\mathbf{H}\in\mathbb{R}^{N\times H}$, slot count $K$,
and iterations $T$:
\begin{enumerate}
    \item Initialise $\mathbf{S}^{(0)} \sim
          \mathcal{N}(\boldsymbol{\mu},\,\mathrm{softplus}(\boldsymbol{\sigma}))$
          from learned parameters.
    \item Compute keys and values once:
          $\mathbf{k} = W_k\,\mathrm{LN}(\mathbf{H})$,
          $\mathbf{v} = W_v\,\mathrm{LN}(\mathbf{H})$.
    \item For $t = 1,\ldots,T$:
    \begin{enumerate}
        \item $\mathbf{q} \leftarrow W_q\,\mathrm{LN}(\mathbf{S}^{(t-1)})$
        \item $\mathbf{L} \leftarrow \mathbf{q}\mathbf{k}^\top / \sqrt{H}
              \in\mathbb{R}^{K\times N}$
        \item $\mathbf{A} \leftarrow \mathrm{softmax}_{K}(\mathbf{L})$
              \hfill (softmax over slots, not over patches)
        \item $\tilde{\mathbf{A}} \leftarrow
              \mathbf{A} \oslash (\mathbf{A}\mathbf{1}_N + \varepsilon)$
              \hfill (normalise each slot)
        \item $\mathbf{U} \leftarrow \tilde{\mathbf{A}}\,\mathbf{v}$
              \hfill (weighted mean of patch values per slot)
        \item $\mathbf{S}^{(t)} \leftarrow
              \mathrm{GRU}(\mathbf{U},\,\mathbf{S}^{(t-1)})$
        \item $\mathbf{S}^{(t)} \leftarrow
              \mathrm{LN}\!\left(\mathbf{S}^{(t)} +
              \mathrm{FFN}(\mathbf{S}^{(t)})\right)$
    \end{enumerate}
    \item Return $\mathbf{S}^{(T)}\in\mathbb{R}^{K\times H}$.
\end{enumerate}

\noindent Default: $K\in\{8,16\}$, $T=3$.

\subsubsection{Cross-Modal Transformer}

Slot tokens from all present modalities are concatenated and refined by an
$L$-layer Transformer encoder:
\begin{equation}
    \mathbf{Z}^{(\ell+1)} = \mathrm{FFN}\!\left(
        \mathrm{LN}\!\left(\mathbf{Z}^{(\ell)} +
            \mathrm{MHA}\!\left(\mathbf{Z}^{(\ell)},
                               \mathbf{Z}^{(\ell)},
                               \mathbf{Z}^{(\ell)}\right)\right)\right),
    \label{eq:xfmr}
\end{equation}
with $L=2$, $n_h=4$ heads. FFN uses a GELU activation with hidden dim $2H$.

\subsubsection{Final Pooling}

Two pooling modes are supported after the cross-modal transformer:

\paragraph{ABMIL pooling.}
Standard gated attention (Eq.~\eqref{eq:abmil}) applied to the slot token
sequence.

\paragraph{CLS pooling.}
A learned CLS token is prepended to the token sequence.
Cross-attention from the CLS token to the remaining tokens is computed via
multi-head attention, followed by layer normalisation.
The CLS output is used as the bag representation.

\subsection{Fusion Variants}

\subsubsection{Variant 1: Early Fusion (\texttt{early} / \texttt{early\_cls})}

All modality patches are mapped to the shared feature space $\mathbb{R}^H$
via their respective frozen backbone, concatenated into a single sequence, and
pooled with ABMIL or a CLS token:
\begin{equation}
    \mathbf{H}_\text{all} = \left[\mathbf{H}_\text{HE}\;;\;\mathbf{H}_\text{BAL}\;;\;
        \mathbf{H}_\text{CT}\;;\;\mathbf{H}_\text{Clin}\right]
    \in\mathbb{R}^{(\sum_m N_m)\times H},
\end{equation}
\begin{equation}
    \hat{y} = f_\text{head}\!\left(\text{pool}(\mathbf{H}_\text{all})\right).
\end{equation}
Modalities absent for a patient are simply omitted.
Patches are sub-sampled to at most $P_\text{max}=2048$ per modality.

\subsubsection{Variant 2: Late Fusion (\texttt{late})}

Each modality is processed independently through its ABMIL encoder to produce
a logit $\hat{y}_m$.
The final prediction is a softmax-weighted sum of present modality logits:
\begin{equation}
    \hat{y} = \sum_{m\in\mathcal{M}} w_m\,\hat{y}_m,\qquad
    \mathbf{w} = \text{softmax}\!\left(\boldsymbol{\ell}[\mathcal{M}]\right),
    \label{eq:late}
\end{equation}
where $\boldsymbol{\ell}\in\mathbb{R}^{|\mathcal{M}_0|}$ is a learned weight
vector over the full modality set $\mathcal{M}_0$, indexed by the present set
$\mathcal{M}$.

\subsubsection{Variant 3: Middle Fusion (\texttt{middle} / \texttt{middle\_cls})}

Each modality is reduced to a single summary vector by ABMIL
(Eqs.~\eqref{eq:abmil}--\eqref{eq:rep}).
The $|\mathcal{M}|$ modality summaries are passed through the cross-modal
transformer (Eq.~\eqref{eq:xfmr}) and pooled:
\begin{equation}
    \hat{y} = f_\text{head}\!\left(
        \text{pool}\!\left(\mathrm{XFmr}\!\left(
            [\mathbf{r}_m]_{m\in\mathcal{M}}
        \right)\right)\right).
    \label{eq:middle}
\end{equation}
The transformer is only applied when $|\mathcal{M}|\geq 2$.

\subsubsection{Variant 4: Cross-Attention Fusion (\texttt{crossattn} / \texttt{crossattn\_cls})}

This variant first applies bidirectional patch-level cross-attention across
modality pairs, then compresses each modality with Slot Attention.
Cross-attention operates between modality patch sequences, not between
individual patch instances across all modalities at once.
For each pair $(m, m')$, MCAT-style~\cite{chen2021mcat} bidirectional
cross-attention produces cross-enriched patch features $\tilde{\mathbf{H}}_m$.
Each enriched sequence is then summarised into $K$ slot vectors by Slot
Attention.
The slot summaries from all modalities are concatenated and fused via the
cross-modal transformer:
\begin{align}
    \tilde{\mathbf{H}}_m &= \text{BidirCA}(\mathbf{H}_m,\, \mathbf{H}_{m'})
        \quad \forall m'\neq m, \\
    \mathbf{S}_m &= \text{SlotAttn}(\tilde{\mathbf{H}}_m)
        \in\mathbb{R}^{K\times H}, \\
    \hat{y} &= f_\text{head}\!\left(
        \text{pool}\!\left(
            \mathrm{XFmr}\!\left(
                [\mathbf{S}_m]_{m\in\mathcal{M}}
            \right)
        \right)\right).
\end{align}
Patch sequences for the bidirectional cross-attention are capped at 256 patches
per modality to limit memory.

\subsubsection{Variant 5: Cross-Modal Slot Fusion (\texttt{crossmodal} / \texttt{crossmodal\_cls})}

Each modality's patch sequence is first summarised into $K$ slot vectors by
Slot Attention.
No patch-level cross-attention between modalities is performed.
The slot summaries serve as compact modality representations.
Cross-modal interaction then happens at the level of these slot summaries
through the cross-modal transformer:
\begin{equation}
    \mathbf{S}_m = \text{SlotAttn}(\mathbf{H}_m), \quad
    \hat{y} = f_\text{head}\!\left(
        \text{pool}\!\left(
            \mathrm{XFmr}\!\left(
                [\mathbf{S}_m]_{m\in\mathcal{M}}
            \right)
        \right)\right).
\end{equation}
All modalities share a single Slot Attention module but retain separate
backbone encoders.

\subsubsection{Variant 6: Iterative Cross-Modal Fusion (\texttt{iterative} / \texttt{iterative\_cls})}

Inspired by~\cite{jaegle2021perceiver,tsai2019multimodal}, we apply $R$
iterative blocks.
Each block consists of within-modal self-attention followed by cross-modal
cross-attention at the patch level.
Let $\mathbf{H}_m^{(r)}$ denote the patch features of modality $m$ after
block $r$.

\paragraph{Self-attention block.}
\begin{equation}
    \tilde{\mathbf{H}}_m^{(r)} = \mathrm{FFN}\!\left(
        \mathrm{LN}\!\left(
            \mathbf{H}_m^{(r)} + \mathrm{MHA}(
                \mathbf{H}_m^{(r)},
                \mathbf{H}_m^{(r)},
                \mathbf{H}_m^{(r)})\right)\right).
\end{equation}

\paragraph{Cross-attention block.}
\begin{equation}
    \mathbf{H}_m^{(r+1)} = \mathrm{FFN}\!\left(
        \mathrm{LN}\!\left(
            \tilde{\mathbf{H}}_m^{(r)} + \mathrm{MHA}\!\left(
                \tilde{\mathbf{H}}_m^{(r)},\,
                \mathbf{C}_{-m}^{(r)},\,
                \mathbf{C}_{-m}^{(r)}
            \right)\right)\right),
    \label{eq:iterative_ca}
\end{equation}
where $\mathbf{C}_{-m}^{(r)} = \bigl[\tilde{\mathbf{H}}_{m'}^{(r)}\bigr]_{m'\neq m}$
is the concatenated patch context from all other modalities.

\paragraph{After $R$ blocks,} each modality is compressed via Slot Attention
and fused with the cross-modal transformer identically to Variant 5:
\begin{equation}
    \hat{y} = f_\text{head}\!\left(
        \text{pool}\!\left(
            \mathrm{XFmr}\!\left(
                [\text{SlotAttn}(\mathbf{H}_m^{(R)})]_{m\in\mathcal{M}}
            \right)
        \right)\right).
\end{equation}
Default: $R=2$, $K\in\{8,16\}$.
HE patches are capped at 1024 inside the iterative blocks to control memory.

\subsection{Phase 2 Summary Table}

\begin{table}[h]
\centering
\caption{Phase 2 fusion variant overview. ``CA'' = cross-attention,
  ``SA'' = slot attention, ``XFmr'' = cross-modal transformer,
  ``pool'' = ABMIL or CLS.}
\label{tab:p2_variants}
\begin{tabular}{llll}
\toprule
Variant & Fusion level & Key components & CLS variant \\
\midrule
\texttt{early}      & patch      & concat $\to$ pool
    & \texttt{early\_cls} \\
\texttt{late}       & logit      & per-mod ABMIL, learned weight
    & {---} \\
\texttt{middle}     & summary    & ABMIL $\to$ XFmr $\to$ pool
    & \texttt{middle\_cls} \\
\texttt{crossattn}  & patch+slot & bidir-CA $\to$ SA $\to$ XFmr $\to$ pool
    & \texttt{crossattn\_cls} \\
\texttt{crossmodal} & slot       & SA $\to$ XFmr $\to$ pool
    & \texttt{crossmodal\_cls} \\
\texttt{iterative}  & patch+slot & $R\times$(SA+CA) $\to$ SA $\to$ XFmr $\to$ pool
    & \texttt{iterative\_cls} \\
\bottomrule
\end{tabular}
\end{table}

\subsection{Phase 2 Hyperparameters}

\begin{table}[h]
\centering
\caption{Phase 2 hyperparameters.}
\label{tab:p2_hparams}
\begin{tabular}{llr}
\toprule
Parameter & Symbol & Value \\
\midrule
Learning rate       & $\eta_2$ & $5\times10^{-5}$ \\
Weight decay        &          & $1\times10^{-3}$ \\
Max epochs          &          & 200 \\
Eval every          &          & 20 epochs \\
Grad accumulation   &          & 4 \\
Transformer layers  & $L$      & 2 \\
Attention heads     & $n_h$    & 4 \\
Attention dropout   &          & 0.1 \\
Modal dropout       &          & 0.3 \\
Max patches (total) & $P_\text{max}$ & 2048 \\
Max HE (iterative)  &          & 1024 \\
Slot iterations     & $T$      & 3 \\
Slot count          & $K$      & 8, 16 \\
Iterative blocks    & $R$      & 2 \\
\bottomrule
\end{tabular}
\end{table}

% =============================================================================
\section{Training Protocol}
% =============================================================================

\paragraph{Optimiser.}
Adam~\cite{kingma2014adam} with the learning rates and weight decay in
Tables~\ref{tab:p1_hparams} and~\ref{tab:p2_hparams}.
Mixed precision (FP16) training is used via PyTorch AMP with gradient scaling.
Gradient accumulation runs over 4 steps.

\paragraph{Cross-validation.}
Nested cross-validation: 5 outer splits $\times$ 4 inner folds.
Model selection is performed on the inner validation fold by balanced accuracy.

\paragraph{Checkpoint selection.}
The checkpoint with the highest validation balanced accuracy is selected as the
best model for each fold.

% =============================================================================
\section{Cross-Attention Interpretability}
% =============================================================================

\subsection{ABMIL Attention Maps}

The per-instance attention weights $\{a_i\}$ from Eq.~\eqref{eq:abmil}
provide a scalar importance score for every patch in each modality.
Patches with high attention are the ones the model relies on most for its
prediction.
For H\&E, these patches can be projected back onto the biopsy slide using the
stored tile coordinates \texttt{tile\_left} and \texttt{tile\_top}, producing
a spatial attention heatmap over the tissue section.
For clinical features, each row of the bag (one feature token) has an
attention weight.
The token label is bin-specific (e.g.\ \texttt{FEV1\_q3} or
\texttt{BMI\_nan}), so the attention map reveals not only which clinical
variables were attended but also whether the attended values were high or low.

\subsection{Clinical $\times$ Cluster Cross-Attention}

After Phase 2 inference, we compute an interpretability map that quantifies
how strongly each clinical feature attends to each cell-cluster type.

Let $\mathbf{F}_\text{Clin}\in\mathbb{R}^{F\times H}$ be the backbone
projections of the $F$ clinical-feature token rows, and let
$\mathbf{C}_m\in\mathbb{R}^{K_m\times H}$ be the slot centroid tokens for
modality $m$.
All cluster tokens are concatenated:
$\mathbf{C} = [\mathbf{C}_m]_{m\in\mathcal{M}}\in\mathbb{R}^{K_\text{tot}\times H}$.

The clinical-cluster cross-attention matrix is:
\begin{equation}
    \mathbf{A}_\text{clin} = \text{softmax}\!\left(
        \frac{\mathbf{F}_\text{Clin}\,\mathbf{C}^\top}{\sqrt{H}}
    \right) \in [0,1]^{F\times K_\text{tot}},
    \label{eq:clin_clust_attn}
\end{equation}
where the softmax is over $K_\text{tot}$.
Each clinical feature produces a probability distribution over all cluster
tokens.

Entry $\mathbf{A}_\text{clin}[f,k]$ quantifies how much clinical feature $f$
attends to cluster $k$.
Each clinical feature token is labelled with its bin-specific identifier
(e.g.\ \texttt{FEV1\_q3}, \texttt{BMI\_nan}), so the attention map
distinguishes which quantile bin of each feature is co-active with which
cell population.
Cohort-level summaries are obtained by averaging $\mathbf{A}_\text{clin}$
over patients separately for rejection-positive and rejection-negative cases
and computing their difference.

% =============================================================================
\section{Cluster-Level Interpretability}
% =============================================================================

Each instance in a bag (patch, cell, or CT slice) carries a cluster-type
annotation string (e.g.\ \texttt{HE\_cells/Lymphocytes}) loaded from the
pre-computed \texttt{.pt} file.
These annotations enable five complementary cluster-level analyses that reveal
which cell populations and clinical states the model attends to.

\subsection{ABMIL Attention by Cluster Type}

For each modality $m$, the ABMIL softmax attention weights
$\boldsymbol{\alpha}_m \in \mathbb{R}^{N_m}$ assign a scalar importance to
every instance.
Given cluster-type labels $c_i \in \mathcal{C}_m$ for instance $i$, the mean
attention for cluster type $c$ is:
\begin{equation}
    \bar{\alpha}_{m,c}
    = \frac{1}{|\mathcal{I}_c|}\sum_{i \in \mathcal{I}_c} \alpha_i,
    \qquad \mathcal{I}_c = \{i : c_i = c\},
    \label{eq:abmil_cluster}
\end{equation}
where $\mathcal{I}_c$ is the index set of all instances of type $c$ in the
bag.

A bar chart of $\bar{\alpha}_{m,c}$ averaged over the cohort, shown separately
for ACR-positive and ACR-negative patients, reveals which cell types are most
attended by the ABMIL head in each clinical group.
Differences between the two groups indicate cell types that are discriminative
for rejection.

\subsection{Modality Importance from Classifier Weights}

For the late-fusion variant, each modality produces an independent logit and
the final decision is a learned weighted sum:
\begin{equation}
    \hat{y} = \sum_{m} w_m \hat{y}_m,
    \qquad \{w_m\} = \text{softmax}(\boldsymbol{\ell}_\text{modal}),
\end{equation}
where $\boldsymbol{\ell}_\text{modal}$ is a learnable parameter vector.
The weight $w_m$ is the classifier's estimate of how much modality $m$
contributes to the final decision.
Comparing the mean weights $\bar{w}_m^{+}$ (ACR) and $\bar{w}_m^{-}$ (non-ACR)
reveals whether one modality dominates in each class.

\subsection{Self-Attention Cluster Co-Attention}

Within the iterative cross-modal model, each block applies a within-modality
self-attention layer with weight matrix
$\mathbf{W}^{(r)} \in [0,1]^{N\times N}$ (patches as queries and keys).
For the final block $R$, the cluster co-attention matrix
$\mathbf{S}_m \in \mathbb{R}^{C\times C}$ measures how strongly patches of
type $c_1$ attend to patches of type $c_2$:
\begin{equation}
    \mathbf{S}_m[c_1, c_2]
    = \frac{1}{|\mathcal{I}_{c_1}||\mathcal{I}_{c_2}|}
      \sum_{i\in\mathcal{I}_{c_1}}\sum_{j\in\mathcal{I}_{c_2}}
      \mathbf{W}^{(R)}[i,j],
    \label{eq:coattn}
\end{equation}
where the patch sequence is capped at $N_\text{cap}=256$ before aggregation.
This reveals intra-modality cell-cell attention patterns that differ between
ACR and non-ACR patients.

\subsection{Slot Attention Cluster Contribution}

The slot attention module assigns each instance to $K$ slots via a competition
softmax.
Let $\mathbf{A}_m \in \mathbb{R}^{K\times N_m}$ be the final-round slot
assignment matrix (post-softmax) for modality $m$.
The $K\times C$ slot cluster contribution matrix is:
\begin{equation}
    \mathbf{B}_m[k,c]
    = \frac{1}{|\mathcal{I}_c|}\sum_{i\in\mathcal{I}_c} A_m[k,i],
    \label{eq:slot_cluster}
\end{equation}
giving the mean assignment weight that slot $k$ places on instances of type
$c$.

We compute both the post-softmax version (Eq.~\eqref{eq:slot_cluster}) and the
pre-softmax (raw logit) version, where
$A^\text{pre}_m[k,i] = \mathbf{q}_k^\top \mathbf{k}_i \cdot \text{scale}$
before the slot competition normalisation.
Side-by-side heatmaps reveal whether slot specialisation is already present in
the raw scores or emerges only after competition.

\subsection{Cross-Modal Slot Connection}

For slot-based fusion variants, the cross-modal transformer operates on the
concatenated slot sequence
$\mathbf{Z} = [\mathbf{Z}_{m_1}; \ldots; \mathbf{Z}_{m_M}]$, yielding a full
attention matrix $\mathbf{X} \in \mathbb{R}^{K_\text{tot}\times K_\text{tot}}$.
We extract the inter-modality block
$\mathbf{X}_{ab} \in \mathbb{R}^{K_a\times K_b}$ for modality pair
$(m_a, m_b)$.

\paragraph{Three-panel connection figure.}
For a given modality pair, we display:
\begin{enumerate}
    \item \textbf{Left}: $\mathbf{B}_{m_a} \in \mathbb{R}^{K_a\times C_a}$,
          showing which cell types load onto each $m_a$ slot.
    \item \textbf{Centre}: $\mathbf{X}_{ab} \in \mathbb{R}^{K_a\times K_b}$,
          showing cross-modal attention between $m_a$ and $m_b$ slot
          summaries.
    \item \textbf{Right}: $\mathbf{B}_{m_b}^\top \in \mathbb{R}^{C_b\times K_b}$,
          showing which $m_b$ cell types each $m_b$ slot represents.
\end{enumerate}
Reading left to right traces the information path:
$m_a$ cell type $\to$ $m_a$ slot $\to$ cross-modal $\to$ $m_b$ slot $\to$
$m_b$ cell type.

\paragraph{Cluster-to-cluster cross-modal attention.}
To summarise the multi-hop path into a single $C_a\times C_b$ matrix, we
compute:
\begin{equation}
    \mathbf{P}_{ab}
    = \mathbf{B}_{m_a}^\top \,\mathbf{X}_{ab}\, \mathbf{B}_{m_b}
    \;\in\; \mathbb{R}^{C_a\times C_b},
    \label{eq:cluster_xmodal}
\end{equation}
where entry $\mathbf{P}_{ab}[c_1,c_2]$ quantifies the effective attention from
cell type $c_1$ in $m_a$ to cell type $c_2$ in $m_b$.
Heatmaps of $\mathbf{P}_{ab}$ for ACR vs.\ non-ACR, and their difference,
reveal which inter-modality cell-type interactions are specific to rejection.

\paragraph{ACR vs.\ Non-ACR comparison.}
All cluster-level interpretability plots are shown in three panels: mean over
ACR-positive patients, mean over ACR-negative patients, and the signed
difference $\Delta = \mu_\text{ACR} - \mu_\text{non-ACR}$.
Positive entries in $\Delta$ indicate mechanisms that are more active in
rejection.
Negative entries indicate mechanisms that are more active in non-rejection.

% =============================================================================
\bibliographystyle{unsrtnat}
\bibliography{references}
% =============================================================================

\end{document}
